\subsection{The main paired campaign and the effort follow-ups}\label{sec:campaign}

\paragraph{Design} The main campaign ran \NSessions\ sessions on \NScenarios\ scripted scenarios (\MinTurns--\MaxTurns\
turns) on both hosts, in paired waves: every arm of a scenario started together from one frozen workspace, each with
its own cache key. Hypotheses, thresholds and the train/test split were preregistered; the estimator was written after
the data. \Cref{fig:forest} gives the confirmatory test-split ratios. The pair reading conditions on quality and was
written after the data; the arm reading uses every pair and agrees in direction.

\begin{figure}[tbp]
\centering
\pgfplotslegendfromname{forestlegend}\par\smallskip
% Confirmatory cost ratios (test split) with 95% scenario-cluster bootstrap CIs. Log x axis.
\begin{tikzpicture}
\begin{axis}[benchmark, xmode=log, log ticks with fixed point,
  xtick={0.5,0.6,0.7,0.8,0.9,1,1.2,1.4,1.6}, xmin=0.45, xmax=1.75,
  scale only axis, width=10.5cm, xmajorgrids, y tick label style={font=\small}, height=5.2cm, ymin=-0.7, ymax=7.0,
  ytick={0,1,2,4,5,6}, yticklabels from table={generated/data/forest.dat}{label},
  xlabel={Cost ratio, routed arm / plain host (log scale; below 1 saves money)},
  legend to name=forestlegend, legend columns=2, legend cell align=left,
  legend style={/tikz/every even column/.append style={column sep=10pt}}]
\draw[black!70, line width=0.8pt] (axis cs:1,-0.7) -- (axis cs:1,6.7);
\draw[okgray, dashed, line width=0.8pt] (axis cs:\HTwoLower,-0.7) -- (axis cs:\HTwoLower,2.6);
\addplot[nodes near coords, point meta=explicit symbolic, nodes near coords style={font=\scriptsize, fill=white, fill opacity=0.85, text opacity=1, inner sep=1pt, anchor=south, yshift=2.5pt, text=black!80}, only marks, mark=*, mark size=2.6pt, color=okblue,
  error bars/.cd, x dir=both, x explicit, error bar style={line width=1pt}]
  table[meta=v, x=pair, y=y, x error minus=pm, x error plus=pp]{generated/data/forest.dat};
\addlegendentry{pair reading (co-primary)}
\addplot[only marks, mark=diamond, mark size=3pt, color=okorange, mark options={line width=1pt},
  error bars/.cd, x dir=both, x explicit, error bar style={line width=0.6pt, densely dashed}]
  table[x=arm, y expr=\thisrow{y}-0.3, x error minus=am, x error plus=ap]{generated/data/forest.dat};
\addlegendentry{arm reading (co-primary, unfiltered)}
\end{axis}
\end{tikzpicture}

\smallskip
\begin{tikzpicture}
\begin{axis}[benchmark, xmode=log, log ticks with fixed point,
  xtick={0.5,0.6,0.7,0.8,0.9,1,1.2,1.4,1.6}, xmin=0.45, xmax=1.75,
  scale only axis, width=10.5cm, xmajorgrids, y tick label style={font=\small}, height=1.6cm, ymin=-0.7, ymax=1.95,
  ytick={0,1}, yticklabels from table={generated/data/forest-h3.dat}{label},
  xlabel={Cost ratio, sticky / shipped on the same scenario (H3; log scale)}]
\draw[black!70, line width=0.8pt] (axis cs:1,-0.7) -- (axis cs:1,1.7);
\draw[okgray, dashed, line width=0.8pt] (axis cs:\HThreeUpper,-0.7) -- (axis cs:\HThreeUpper,1.7);
\addplot[nodes near coords, point meta=explicit symbolic, nodes near coords style={font=\scriptsize, fill=white, fill opacity=0.85, text opacity=1, inner sep=1pt, anchor=south, yshift=2.5pt, text=black!80}, only marks, mark=*, mark size=2.6pt, color=okblue,
  error bars/.cd, x dir=both, x explicit, error bar style={line width=1pt}]
  table[meta=v, x=pair, y=y, x error minus=pm, x error plus=pp]{generated/data/forest-h3.dat};
\addplot[only marks, mark=diamond, mark size=3pt, color=okorange, mark options={line width=1pt},
  error bars/.cd, x dir=both, x explicit, error bar style={line width=0.6pt, densely dashed}]
  table[x=arm, y expr=\thisrow{y}-0.3, x error minus=am, x error plus=ap]{generated/data/forest-h3.dat};
\end{axis}
\end{tikzpicture}

\caption{On the main-v1 test split, every Fable routing arm saved money and every Opus arm cost more. Confirmatory cost ratios on the main-v1 test split. Top: each routed arm against the plain host. Bottom: the
sticky arm against the shipped arm on the same scenarios. Whiskers are 95\,\% intervals; the dashed lines mark
preregistered thresholds.}
\label{fig:forest}
\howtoread{The dot is the geometric-mean cost ratio, the whisker its 95\,\% interval, on a log axis. Rows left of 1
saved money. All Fable rows are well left of 1. All Opus rows are right of it.}
\end{figure}

\paragraph{Effort follow-ups} Two follow-ups measured reasoning effort directly, each on \EcScen\ test scenarios in
concurrent pairs: plain Sonnet at medium cost \EcRatio\X\ its default (\EcCI; designed in advance, with a provenance
limitation recorded in its package) and plain Fable at medium cost \EfRatio\X\ its default (\EfCI; preregistered),
both with no detectable turn-pass loss. On Opus, medium effort did not demonstrate a saving in S1
(\SHFourRatio\X, \SHFourCI).

\paragraph{Pricing every decide-once policy on the existing sessions} Because a decide-once policy's session equals
either the plain-host session or the routed session of the same scenario, every such policy can be priced from the
main campaign without new runs (\cref{fig:a0,tab:a0}) \expl. On Fable, Jev routed \AoFableJevRoute\,\% of sessions
and cost \AoFableJev\X\ the plain host. The simple rule R* (route unless the workspace has more than \AoScopeFiles\
files) routed \AoFableRstarRoute\,\% and cost \AoFableRstar\X\ at the same turn-pass difference (\AoFableRstarDtp).
Jev's ``keep on host'' calls bought nothing: in every scenario-rep where Jev kept the host, the routed session scored
the same. Jev and R* disagreed on only \AoJevRstarN\ of \AoJevRstarUnits\ scenario-reps, which bounds how much they can
differ. On Opus the price gate, which never routes there, sat at the oracle (regret \$\AoOpusShippedRegret\ per 1,000
sessions), while ungated Jev would have cost \AoOpusJev\X. The same re-analysis found the cache-rewrite cost of effort
switches shown in the main paper (\cref{sec:dtrace}).

\begin{figure}[tbp]
\centering
\pgfplotslegendfromname{aolegend}\par\smallskip
% A0 counterfactual: decide-once policies on main-v1 (all 70 scenarios x 2 reps), geometric-mean cost ratio vs the
% plain host with 95% scenario-cluster bootstrap CIs. Log x axis.
\begin{tikzpicture}
\begin{axis}[benchmark, xmode=log, log ticks with fixed point, xtick={0.5,0.6,0.8,1,1.2}, xmin=0.45, xmax=1.35,
  scale only axis, width=8.0cm, height=5.2cm, xmajorgrids, y tick label style={font=\small},
  ymin=-0.7, ymax=8.0, ytick={\AoTicks}, yticklabels/.expanded={\AoLabels},
  xlabel={Cost ratio, policy / plain host (log scale; below 1 saves money)},
  title={Who decides, and what it costs}, title style={font=\small\bfseries},
  legend to name=aolegend, legend columns=1, legend cell align=left]
\draw[black!70, line width=0.8pt] (axis cs:1,-0.7) -- (axis cs:1,7.7);
\addplot[nodes near coords, point meta=explicit symbolic, nodes near coords style={font=\scriptsize, fill=white, fill opacity=0.85, text opacity=1, inner sep=1pt, anchor=south, yshift=2.5pt, text=black!80}, only marks, mark=*, mark size=2.6pt, color=okblue, error bars/.cd, x dir=both, x explicit,
  error bar style={line width=1pt}] table[meta=v, x=r, y=y, x error minus=m, x error plus=p]{generated/data/a0-policies.dat};
\addlegendentry{policy on main-v1, all scenarios (95\,\% CI)}
\end{axis}
\end{tikzpicture}

\caption{On Fable the rule R* is at least as cheap as Jev; on Opus every routing policy costs more. Decide-once policies priced on the existing main-v1 sessions (all scenarios, both reps), with 95\,\%
scenario-cluster intervals. The oracle chooses with hindsight on the other rep, at equal quality. Exploratory.}
\label{fig:a0}
\howtoread{Each row is a policy. Left of 1 saves money. On Fable, the rule R* is at least as cheap as Jev. On Opus,
every policy that routes costs more than keeping the host.}
\end{figure}

\begin{table}[tbp]
\centering\scriptsize
\caption{Decide-once policies on main-v1 (exploratory): share routed, geometric-mean cost ratio and turn-pass
difference against the plain host, and dollars per 1,000 sessions.}
\label{tab:a0}
\begin{tabular*}{\textwidth}{@{\extracolsep{\fill}}l r r r r@{}}
\toprule
Policy (decided once per session) & routed & cost ratio [95\,\% CI] & turn-pass $\Delta$ [95\,\% CI] & \$ / 1,000 \\
\midrule
Fable: always host (plain) & 0\,\% & 1.000 [1.000, 1.000] & 0.000 [0.000, 0.000] & 5{,}311 \\
Fable: Jev decides once & 89\,\% & 0.545 [0.510, 0.585] & $-$0.019 [$-$0.052, 0.010] & 3{,}108 \\
Fable: rule R* (always route, scope gate) & 99\,\% & 0.509 [0.484, 0.538] & $-$0.019 [$-$0.052, 0.010] & 2{,}739 \\
Fable: oracle, equal quality & 81\,\% & 0.583 [0.544, 0.627] & 0.012 [$-$0.004, 0.030] & 3{,}268 \\
Opus: shipped (price gate, never routes) & 0\,\% & 1.000 [1.000, 1.000] & 0.000 [0.000, 0.000] & 2{,}125 \\
Opus: Jev decides once, no gate & 89\,\% & 1.159 [1.092, 1.232] & $-$0.009 [$-$0.033, 0.013] & 2{,}498 \\
Opus: always route & 100\,\% & 1.197 [1.125, 1.273] & $-$0.005 [$-$0.030, 0.019] & 2{,}620 \\
Opus: oracle, equal quality & 22\,\% & 0.991 [0.975, 1.008] & 0.000 [$-$0.011, 0.009] & 2{,}124 \\
\bottomrule
\end{tabular*}

\end{table}

